\begin{helper}
Let $\mu$ be the push-forward law of the activation--priority output set in
\noderef{RandomIndependentSetSampling}, and let $\nu$ be the underlying law on
pairs $(A,\pi)$.  For every finite set of vertices $X$,
\[
\operatorname{ofReal}\!\left(\sum_S
  {\bf 1}_{S\cap X\ne\emptyset}\,\mu(S)\right)
=
\nu\!\left(\{(A,\pi):I(A,\pi)\cap X\ne\emptyset\}\right),
\]
where
\[
I(A,\pi)=\{z\in A:\pi(w)<\pi(z)\text{ for every }w\in A\cap N_G(z)\}.
\]
\end{helper}
\begin{proof}
Let \(I(A,\pi)\) denote the output set appearing in
\noderef{RandomIndependentSetSampling}.  Put
\[
\mathcal T=\{S:S\text{ is a finite vertex set and }S\cap X\ne\emptyset\}.
\]
Since the weights \(\mu(S)\) are nonnegative, finite additivity of
\(\operatorname{ofReal}\) over nonnegative real sums rewrites
\[
\operatorname{ofReal}\!\left(\sum_S
  {\bf 1}_{S\cap X\ne\emptyset}\mu(S)\right)
\]
as
\[
\sum_{S\in\mathcal T}\operatorname{ofReal}(\mu(S)).
\]
For each finite set \(S\), the push-forward clause of
\noderef{RandomIndependentSetSampling} identifies
\(\operatorname{ofReal}(\mu(S))\) with
\[
\nu\{(A,\pi):S=I(A,\pi)\}.
\]
The events \(\{S=I(A,\pi)\}\) are pairwise disjoint as \(S\) ranges over
distinct finite sets.  Their union over \(S\in\mathcal T\) is exactly the
event \(\{(A,\pi):I(A,\pi)\cap X\ne\emptyset\}\): if the output meets \(X\),
choose \(S=I(A,\pi)\), and conversely if \(S=I(A,\pi)\) and \(S\cap X\ne
\emptyset\), then \(I(A,\pi)\cap X\ne\emptyset\).  Finite additivity of
\(\nu\) over this disjoint union gives the displayed identity.
\end{proof}
